\documentclass[t,aspectratio=169]{beamer}

\usepackage{soul}
\usepackage{tikz}
\usepackage{multicol}
%\usepackage{gensymb}
\usepackage[utf8]{inputenc}

% ----------------------------------------------------------------------------------------------------------------------------------------

\title{\includegraphics[width=45pt]{oepecsetlogo.png}\\(KTXFI1EBNF) Physics I. Practice}

\author{Dr.~Gábor FACSKÓ, PhD\\\tiny Assistant Professor / Senior Research Fellow\\\textit{facsko.gabor@uni-obuda.hu}}

\institute{\Tiny Óbuda University, Faculty of Electrical Engineering, 1084 Budapest, Tavaszmező u. 17.\\Wigner Research Centre for Physics, Department of Space Physics and Space Technology, 1121 Budapest, Konkoly-Thege Miklós út 29–33.\\ \url{https://wigner.hu/~facsko.gabor}}

\date{March~9,~2026}

\logo{\includegraphics[height=0.9 true cm]{kekcsik.png}}

% ----------------------------------------------------------------------------------------------------------------------------------------

\begin{document}

\frame{\titlepage}

% ----------------------------------------------------------------------------------------------------------------------------------------

\begin{frame}[fragile,squeeze,allowframebreaks]
\frametitle{Exercises}
\begin{itemize}
\setlength{\parskip}{0pt}
\setlength{\itemsep}{0pt}
\item \textbf{Thermodynamics I:} Temperature scales, thermal expansion, thermometers. State variables, ideal gas, state changes, ideal gas equation of state, Boyle’s law, Gay-Lussac’s laws, universal (combined) gas law, heat, specific heat, molar heat, heat capacity.


\pagebreak
\item Problems:
\begin{enumerate}
\item Convert the $68{}^{\circ}F$ temperature value to Celsius and Kelvin.

\item An aluminium 100\,km long transmission line was installed at $0{}^{\circ}C$. What will be its length in $40{}^{\circ}C$ during summer? The thermal coefficient of the aluminium material is $2.4 \cdot 10^{-5}\frac{1}{K}$.

\item How much heat energy do we need to transfer to 3\,kg of water to increase its temperature from $30{}^{\circ}C$ to $50{}^{\circ}C$? The specific heat capacity is $4.19 \cdot 10^{3}\frac{J}{kg\,K}$.
  
\item We place a copper piece with a mass of 40 gramms and a temperature of $20{}^{\circ}C$ into 200 gramms of water at $80{}^{\circ}C$. What will be the common temperature? The specific heat capacity of copper is $3.85 \cdot 10^{2}\frac{J}{kg\,K}$ and the specific heat capacity of water is $4.19 \cdot 10^{3}\frac{J}{kg\,K}$.

\pagebreak
\item At a temperature of $20{}^{\circ}C$, a steel shaft with a diameter of $11.28\,cm$ needs to be fitted with an aluminum ring with an inner diameter of $11.25\,cm$ at the same temperature. The linear thermal expansion coefficient of steel is $12 \cdot 10^{-6}\frac{1}{K}$, and the linear thermal expansion coefficient of alumium is $28.7 \cdot 10^{-6}\frac{1}{K}$. What temperature should the ring be heated to?

\item The atmospheric balloon floating at an altitude of several kilometers has a volume of $800\,m^{3}$. The temperature and pressure of the helium inside it are the same as the surrounding $-50{}^{\circ}C$ and $5.26 \cdot 10^{3}\,Pa$, respectively.
\begin{enumerate}
\item[(a)] How many moles of helium are in the balloon?
\item[(b)] What is the mass of the helium gas? $\left(M = 4 \frac{g}{mol}\right)$
\item[(c)] What was the volume of the balloon at the time of launch on Earth, if the gas inside the balloon was characterized by $0{}^{\circ}C$ and $10^{5}\,Pa$ pressure?
\end{enumerate}
\end{enumerate}

\end{itemize}

\end{frame}

% ----------------------------------------------------------------------------------------------------------------------------------------

\begin{frame}
%\frametitle{Minden jó, ha a vége jó}

\vfill
\begin{center}
\textbf{\huge The End}

\bigskip
Thank you for your attention!
\end{center}
\vfill
%\textbf{Acknowledgements} Thank you for the course materials for Szilárd~ZSÓKA and Zoltán~VARGA.
\end{frame}

\end{document}
